\chapter{Materials, procedure and availability}

\section{Source texts}

The papers are the published proceedings of the International Conference on Salish and Neighbouring
Languages, distributed as PDF files by the University of British Columbia Working Papers in
Linguistics. The archive holds 993 papers; the automatic extractors need eleven of them, and the
screening reads every paper with no extractor.

Two conversion choices have a wrong answer that looks like a right one.

\textbf{The encoding.} A PDF-to-text converter that writes Latin-1 has neither ʔ, ə nor ɬ. A paper
converted that way still opens and still looks like a paper, with the language taken out of it.
Every paper here is converted to UTF-8.

\textbf{The reading order.} Converted in layout order, a running header prints above the title it
sits under and a two-column table interleaves. The papers are read in the order the PDF stores the
page, and the automatic extractors are written against that order.

A converted paper whose fonts renumber their glyph codes and declare no map back to Unicode is
marked as not representing the page, with the fonts named. Its text is kept and nothing downstream
reads it as the paper. Six of the 146 PDFs held are in that class, and the references chapter
names them.

\section{When the text is not the page}

Two of the eleven papers, Lyon \cite{src:Lyon-2015} and Lindley and Lyon \cite{src:Lindley-and-Lyon-2013},
are set in a font of that class: their text layer holds the font's alphabet and not what the page
prints. For those two each page is rendered as an image at three times its size, which puts a 12 pt
body around 50 pixels tall, where a stacked diacritic stops being a guess. A page text is drafted
from the font's codes, line for line with the extraction, and graded against the hand extraction.
Two questions cannot be settled without looking at the page: which \texttt{w} is a labialized
consonant, and which inserted space is a word boundary.

The drafted page text replaces the font repair for those papers. Applying the mapping to text it has
already been applied to destroys the text.

\section{Procedure}

\begin{enumerate}
  \item Convert each paper to text, as above.
  \item For the two papers whose text is not the page, render the pages and draft a page text.
  \item Run each paper's automatic extractor, which writes that paper's gold standard lines, its
        other text, and the lines it could not sort.
  \item Check each hand extraction against its paper in both directions. A form written down that
        the paper does not hold is a typing slip. A word in the paper that no row holds is a row
        somebody skipped. A person transcribing a thirty-seven page paper into a table produces
        both. For the two drafted papers, the first direction measures the draft and not the table.
  \item Check each automatic extractor against the hand extraction: whether it found each form,
        whether it assigns the same kind, and whether it assigns the same dialect. Kind decides the
        corpus. A rejected tableau candidate filed as a citation form puts a form the paper's own
        analysis rejects into the gold standard corpus.
  \item Check coverage: every language token in each paper against the corpus built from it, both
        sides through the same repairs. For the two drafted papers a coverage of zero missing means
        the extractor got everything the font's encoding held, which is weaker.
  \item Calibrate the screening against the gold standard corpora and the English spans the hand
        extractions mark, then screen every paper without an extractor against both profiles.
  \item Read which language each paper is about out of its own title and abstract, ask whether the
        distributions agree, and admit candidates to each corpus in batches while it stays on its
        own growth curve.
\end{enumerate}

Two further tests support the repairs. A candidate character mapping is tested by counting how many
damaged tokens become forms attested in a clean paper; the reading is the change between the two
rates. The screened output is matched case sensitively against the gold standard corpus, and a form
that matches only once case is folded is the same word written wrong; the count of those is the
formatting damage.

\section{Shared definitions}

One definition, the text space, fixes which characters these orthographies are written with, and
every component reads it. A per-paper set is the text space plus what that paper adds. The
repairs are each one function: putting back words split internally by an extraction, closing the
space a PDF leaves after a stacked diacritic (with the mark set each paper's spacing allows), reading
a paper's glyph names back as the characters they name, and rejoining words broken across lines.

\section{The sound representation}

Each recording is reduced to one row per 10 ms frame: a 24 bit segment field and a 4 bit prosody
field. The segment field is the loudness heard in log spaced bands, made shift invariant and rotated
onto the recording's own singular vectors; the prosody field is loudness, pitch against the
speaker's median, pitch movement and voicing. The delta from a flat distribution is read at 6, 8,
10, 12, 14 and 24 bits. The narrow widths are the ones to read: a 24 bit field has 16.8 million
states and the longest recording has 135 thousand frames.

\section{Data and code availability}

The source papers and recordings are published by UBCWPL and are available from its archive. The
hand extractions, the automatic extractors, the checks, the screening and the sound representation
are in the orior repository, \texttt{https://github.com/dstroy0/orior}. Every number in this research paper
is computed by that code from those inputs.
